%=== CHAPTER TWO ===
\chapter{Background and Related Work}
\label{ch:literature}
\begin{spacing}{1.5}
\setlength{\parskip}{0.22in}

\section{Model Predictive Control}

MPC repeatedly solves a finite-horizon optimal control problem, applies the first input and shifts the horizon forward. Its practical value comes from combining prediction, multivariable coordination and explicit constraints within one optimisation \cite{rawlings2017mpc,mayne2014mpc}. Linear MPC is mature and computationally efficient, but a single linear model is poorly suited to the large attitude changes and thrust-direction coupling of a quadrotor. NMPC retains the nonlinear dynamics and nonlinear tracking errors directly \cite{grune2017nmpc}.

The price is computation. Multiple shooting, sequential quadratic programming (SQP), Gauss--Newton Hessian approximations and real-time iteration schemes are widely used to make NMPC practical \cite{diehl2005rti}. CasADi provides symbolic differentiation and a high-level optimisation interface \cite{andersson2019casadi}; acados generates efficient optimal-control solvers around tailored SQP and QP components \cite{verschueren2022acados}. The present work uses full SQP with merit-function backtracking because preliminary figure-eight tests exposed convergence failures under a fixed full step.

\section{Quadrotor NMPC}

A quadrotor is differentially flat under common assumptions, and many successful controllers exploit position/yaw references and a cascaded inner attitude loop. NMPC can instead optimise collective thrust and body torque against the complete nonlinear rigid-body model. Comparisons on rotary-wing platforms have shown why the nonlinear formulation becomes advantageous as trajectory curvature and operating envelope increase \cite{kamel2017linear}.

The controller in this dissertation predicts position, velocity, unit quaternion and body rate. The attitude residual is the vector part of the relative quaternion, which avoids treating the four quaternion components as independent Euclidean coordinates. Input limits represent collective thrust and body-torque authority. This formulation is nominally standard; the distinguishing issue is that the equilibrium input and the plant parameters change at payload attachment.

\section{Uncertainty and Adaptation in MPC}

Robust MPC protects against bounded uncertainty, while stochastic MPC describes uncertainty probabilistically. Both can preserve margins without explicitly identifying the current plant, but conservatism grows when the uncertainty set must cover both an empty and a loaded vehicle. Adaptive MPC instead updates a parameterised model online. Certainty-equivalent schemes use the current estimate directly, set-membership schemes maintain a feasible parameter set, and dual-control schemes intentionally excite the plant to learn \cite{aswani2013learning,hewing2020learning}.

Learning-based MPC extends this idea using Gaussian processes, neural networks or learned residual models. Data-driven MPC for quadrotors can compensate complex aerodynamic effects \cite{torrente2021datadriven}, and neural approximations can reduce online computation \cite{salzmann2023neural}. Those methods are valuable when the uncertainty is poorly structured. Payload pickup has stronger prior structure: rigid-body mechanics already specify how mass and attachment geometry enter the equations. The present work therefore estimates a small physical parameter set rather than learning a general residual map.

\section{State and Parameter Estimation}

Recursive filters such as the extended Kalman filter propagate a local Gaussian approximation. They are efficient and central to flight-state estimation, but inequality constraints are not naturally hard constraints and abrupt parameter changes can violate the local linearisation assumptions. MHE estimates a window of past states and disturbances by solving a constrained optimisation. It can impose bounds such as a physically admissible mass interval and can attribute residuals using the same nonlinear model as the controller \cite{rawlings2017mpc,rao2003mhe}.

A moving window introduces an important trade-off. A longer horizon improves noise rejection and parameter observability, but delays adaptation after a step change. A shorter horizon responds rapidly but may let measurement noise and transient accelerations contaminate the estimate. The implemented estimator uses 20 stages at \SI{0.1}{s}, giving a \SI{2.0}{s} window aligned with the NMPC/control data stream.

\section{Observability of Payload Mass and Geometry}

Near hover, translational acceleration depends on the ratio between physical thrust and mass. If the estimator is supplied with a thrust command that the NMPC computed using its own current mass estimate, the apparent relation is self-consistent for any estimate and contains no independent information. Rotor speed provides a separate measurement channel through the quadratic thrust map. This distinction between \emph{intended} and \emph{physical} thrust is central to the proposed design.

Mass, inertia and horizontal CoM offset are not equally observable. Mass is excited by ordinary vertical support. Inertia requires angular acceleration, which is weak during much of a delivery mission. A horizontal CoM offset can be recovered from the steady rotor-force imbalance that balances its trim moment. The architecture therefore assigns each parameter to the strongest physical channel rather than augmenting every possible quantity into the MHE state.

\section{Research Gap}

The literature contains strong results on quadrotor NMPC, learned residual dynamics and payload-aware control. The remaining engineering gap addressed here is the complete transition through grasp and release: independent physical-input reconstruction, bounded online mass estimation, propagation of the estimate into the coupled inertia/CoM model, recentering of the NMPC input penalty, coordination with PX4 inner-loop gains and evaluation without feeding payload ground truth to the adaptive controller.

The dissertation consequently treats MHE and NMPC as one architecture rather than two isolated algorithms. Estimation is judged by whether its outputs are identifiable and useful for control; control is judged by whether it remains consistent with the quantities being estimated.

\end{spacing}
\newpage
